\documentclass[10pt,twocolumn]{article}
\usepackage[scaled]{helvet}
\renewcommand\familydefault{\sfdefault}
\usepackage[T1]{fontenc}
\usepackage[margin=0.7in]{geometry}
\usepackage{titlesec}
\usepackage{enumitem}
\usepackage{parskip}
\setlength{\columnsep}{0.24in}
\setlist[itemize]{leftmargin=1.2em, topsep=0.2em, itemsep=0.18em}
\titleformat{\section}{\large\bfseries}{\thesection}{1em}{}

\begin{document}
\title{\vspace{-1.6cm}AWS ECS (Elastic Container Service)}
\author{AWS ML Study Notes}
\date{}
\maketitle
\small

\section*{Overview}
ECS is AWS managed container orchestration. It schedules containers onto compute capacity and keeps services running according to the task definition you declare.

For ML applications, ECS is a practical choice for serving models, running supporting APIs, and executing recurring preprocessing or postprocessing jobs without managing Kubernetes.

\section*{Core Concepts}
\begin{itemize}
\item Core objects include clusters, task definitions, tasks, and services. A task definition describes the containers, CPU, memory, networking, secrets, and startup behavior.
\item ECS can run on EC2 instances that you manage or on Fargate, where AWS manages the underlying servers.
\item Services maintain a desired number of healthy tasks and integrate with load balancers, service discovery, and deployment strategies.
\end{itemize}

\section*{How It Works}
\begin{itemize}
\item You package the application as a container, store it in ECR, create a task definition, and then run it as a task or a service.
\item For a long running inference API, ECS typically deploys tasks behind an Application Load Balancer and scales them based on CPU, memory, request count, or custom metrics.
\item For scheduled or event driven jobs, ECS can launch short lived tasks that process a queue, transform data, or prepare artifacts for later training or inference stages.
\end{itemize}

\section*{AI and ML Relevance}
\begin{itemize}
\item ECS works well for FastAPI model servers, embedding services, feature preprocessing workers, and internal tools that need container consistency but not the full weight of a larger platform.
\item If your inference stack depends on custom native libraries or multiple sidecars, ECS often gives a cleaner operational model than forcing everything into a serverless boundary.
\end{itemize}

\section*{Operational Notes}
\begin{itemize}
\item Right size task resources and container concurrency. Overcommitting CPU or memory leads to noisy failures and bad latency under load.
\item Use IAM task roles, central logging, health checks, and rolling deployments so the service can update without downtime.
\item Fargate reduces operational work, but EC2 backed ECS may be cheaper or more flexible for heavy GPU or host level customization requirements.
\end{itemize}

\section*{What To Remember}
\begin{itemize}
\item Task definition is the runtime contract. Service is the controller that keeps tasks alive.
\item Fargate optimizes convenience; EC2 launch type optimizes control.
\item ECS is a strong middle ground when you want containers without adopting Kubernetes complexity.
\end{itemize}

\end{document}
